\chapter{Introduction}
\label{chap:introduction}

This chapter introduces Medorama: what it does, the programs it is made of, how it protects
patients' identities, the synthetic data used in this manual, and how the manual is organized.

Medorama is a unified patient-journey research platform linking EHR data, physiological
waveforms, imaging and laboratory results across databases, free for noncommercial use. It
brings together what is known about a patient from many sources: admissions and bed stays, laboratory results, vital
signs, medications, conditions, notes, ECGs, imaging and bedside-monitor recordings. Every
patient has a page with their \emph{Patient Journey}, a picture of the whole history with one
column per clinical category, beside a searchable report of every entry. Researchers build
\emph{cohorts} across every database, with criteria written as expressions such as
\code{QTc > 470ms AND female}. Recordings, ECGs, echo loops and images open in Medorama's own
viewers, as tabs of the same page.

\emph{Privacy Mode}, the default for every user except administrators with the PHI permission,
shows the same data without any identifier: the server replaces names and numbers with
placeholders and pseudonyms, moves dates back, and de-identifies files before they leave it.

\screenshot{patient.png}{A patient's page in the dark theme: the Patient Journey (center), the
report and the detail of a 12-lead ECG (right). Synthetic data.}{fig:introduction-patient}

\begin{note}[Research use only]
Medorama is research software, not a medical device: it is not cleared for diagnosis or
treatment, and must not be used to make clinical decisions.
\end{note}

\section{What Medorama does}
\label{sec:introduction-what}

\Cref{tab:introduction-features} summarizes what Medorama offers; each line is a chapter of
\cref{part:using} or \cref{part:data}.

\begin{table}[htbp]
\caption{What Medorama does.}
\label{tab:introduction-features}
\begin{tabularx}{\linewidth}{@{}P{2.9cm}L@{}}
\toprule
{\tabhead Feature} & {\tabhead What it is} \\
\midrule
Cohorts & Patients found across every database by criteria (results, vital signs, ECG and
  imaging measurements, diagnoses, medications, encounters and ICU stays, age, sex, text),
  written as expressions or chosen line by line; kept, refreshed, combined, cloned and shared;
  each cohort's definition records its criteria and every patient added manually; members and
  their data exported as CSV and Excel (\cref{chap:cohorts}). \\
Patient Catalog & Every patient, searchable; each one's page: identity and allergies,
  encounters, medications and conditions, the Patient Journey, a report of every entry
  searchable down to the notes, the detail of each entry (\cref{chap:patients}). \\
Databases & Every imported database as a table, as it was imported, each row linked to its
  patient and its file to its viewer; connected databases (an OMOP CDM instance, a research
  warehouse) read in place, their persons linked to the same patients
  (\cref{chap:databases-tab}). \\
Viewers & VitalStream for ICU monitoring, Holter, EEG and sleep recordings, from weeks down
  to single samples; the ECG viewer, on standard paper; the PACS viewer, for DICOM studies (echo and
  angiography loops, CT, MR, nuclear medicine, X-ray, structured reports)
  (\cref{chap:viewers-in-medorama}). \\
Privacy Mode & Pseudonyms and placeholders instead of identifiers, dates moved back with
  intervals kept, de-identified WFDB, MUSE, DICOM and EDF headers; enforced by the server
  (\cref{sec:introduction-privacy}, \cref{chap:protecting}). \\
Sign-in & Medorama's own accounts and passwords, or the institution's sign-in (single sign-on,
  multi-factor authentication) through an authenticating proxy (\cref{chap:sign-in}). \\
Work & Background indexing, a Holter's heart rate, plug-in actions on any selection of
  recordings (\cref{chap:actions}); users, roles and a separate PHI permission; a sealed
  audit log (\cref{chap:administration}). \\
Scale & SQLite on one server, or PostgreSQL for several servers and many users, with the
  same program (\cref{chap:deployment}). \\
\bottomrule
\end{tabularx}
\end{table}

\section{The parts}
\label{sec:introduction-parts}

Medorama is made of separate programs, each in its own folder of the repository, joined by
written contracts: a file format, an HTTP API, a command-line protocol, an address or a folder
layout. \Cref{tab:introduction-parts} lists them, and \cref{chap:architecture} describes how they
fit together, where each one runs and which trust boundaries separate them.

\begin{table}[htbp]
\caption{The parts of Medorama.}
\label{tab:introduction-parts}
\begin{tabularx}{\linewidth}{@{}P{2.4cm}P{3.6cm}L@{}}
\toprule
{\tabhead Part} & {\tabhead Folder} & {\tabhead What it is} \\
\midrule
Medorama & \file{medorama/} & The server and its web page (Rust, SQLite or PostgreSQL, plain
  JavaScript): the catalog, connected databases, cohorts, the patient pages, the privacy gate,
  users and audit, jobs and plug-ins. Run with \cmd{medorama serve}. \\
VitalStream & \file{viewers/vitalstream/} & The waveform viewer, in three implementations
  (Python, Rust and C++: \cmd{vitalstream-py}, \cmd{vitalstream-rs}, \cmd{vitalstream-cpp}),
  with the same page, the same API and the same index, checked against each other. It reads
  WFDB records, multi-segment records, folders of records, WFDB collections, and EDF, EDF+, BDF
  and BDF+ files, and other formats through adapters. \\
ECG viewer & \file{viewers/ecg/} & 12-lead ECGs on standard paper with their report: WFDB, GE
  MUSE XML, DICOM ECGs, EDF, JSON, PNG, JPEG, PDF. A static page; \cmd{muse-py} reads GE MUSE XML
  and builds ECG databases. \\
PACS viewer & \file{viewers/pacs/} & DICOM studies: echo and angiography loops, CT, MR, NM and
  X-ray series, structured reports, PDFs, pictures and videos, with its own JPEG decoders for
  what browsers do not decode. A static page; \cmd{dicom-py} builds databases of DICOM
  studies. \\
Libraries & \file{libs/} & \code{signal-core} (signal types and values shared by the Rust
  code) and \code{engine} (the crate \code{recording-engine}: indexing and reading
  recordings, and the adapter protocol). WFDB records are read with wfdb-x
  (\url{https://github.com/alphanumericslab/wfdb-x}) in Rust, C++ and Python. \\
Add-ons & Outside the repository & Programs that read formats under a vendor's license (for
  example BedMasterEx \file{.Stp} files, read by stp-adapter, which is private and available to
  institutions under agreement), installed by an administrator with
  \cmd{medorama addons install PROGRAM} and speaking a small open protocol
  (\cref{chap:addons}). \\
\bottomrule
\end{tabularx}
\end{table}

The repository reads open formats only (WFDB, EDF, DICOM, MUSE XML, tables). Nothing licensed is
compiled into Medorama or a viewer: an add-on is a program of its own, distributed under its
own terms.

\paragraph{Each viewer is usable alone.} VitalStream, the ECG viewer and the PACS viewer do not
depend on Medorama. Each installs and runs on its own, and opens the files of one computer:
VitalStream under a root folder fixed when it is built, with an optional password; the ECG and
PACS viewers as pages opened from disk, from any web server, or from the website. Medorama
opens the same viewers, unchanged, at the file to show, inside its own page; patients,
databases, sign-in, users and de-identification belong to Medorama alone.
\Cref{part:viewers} describes the viewers on their own (\cref{chap:viewers-alone}).

\paragraph{Light and dark.} The page and every viewer have a light and a dark theme, and follow
the system's setting until the user chooses one (\cref{chap:page}).

\section{Privacy in one page}
\label{sec:introduction-privacy}

Medorama runs on servers that the institution controls. It listens on the local machine
(127.0.0.1) by default and sends nothing anywhere else. Data enter only as databases that an
administrator imports or connects (connected databases are read in place, read-only), and the
files they name (recordings, ECGs, images) are read where they lie and never written. This
section summarizes the protection of patients' identities; \cref{chap:protecting} explains it in
full.

\paragraph{Identifiers stay in the catalog.} MRNs, names, dates of birth, account numbers and
the imported patient and encounter IDs exist only in Medorama's database. They are shown and
searchable only for users with the \emph{PHI permission} who switched Privacy Mode off, and
those users' searches and views are written to the audit log. The viewers, the plug-ins, the
indexes and the results folders receive pseudonyms only.

\paragraph{Privacy Mode.} A session is in Privacy Mode unless its user has the PHI permission
and has switched it off with the \ui{Privacy Mode} switch in the page's header. Administrators
with the PHI permission start with Privacy Mode off, everyone else with it on, and users without
the permission cannot leave it. In Privacy Mode:
\begin{itemize}
\item patients and encounters are named by pseudonyms (\code{MO-XXXX-XXXX},
  \code{MO-E-XXXX-XXXX}), random and stored, never derived from an identifier;
\item names, MRNs, account numbers and dates of birth become placeholders
  (\code{[Patient Name]}, \code{[MRN]}, \code{[CSN]}, \code{[Date of Birth]}); staff names
  become \code{[Staff Name]};
\item every date moves back by the catalog's offset and by each patient's own number of days,
  drawn once and kept, so that the intervals between a patient's events are kept; ages above
  89 become \code{90+};
\item files are de-identified on the server: WFDB headers, GE MUSE XML and EDF headers are
  rewritten, DICOM headers are rewritten with the pixels untouched, and file names are replaced
  by tokens. Pictures, videos and PDFs are sent as they are, marked with an amber warning
  triangle, because they may show burned-in text. Anything else is refused;
\item the text of notes and reports is shown with the patient's own names and date of birth
  replaced, and marked, because it may still name the patient or others. It is withheld
  (\code{[Withheld]}) when an administrator withholds \ui{The text of notes and reports};
\item searches of identifiers are refused: patients are found by their pseudonyms;
\item exports pass the same gate (their file names end in \code{-privacy}).
\end{itemize}

\paragraph{The gate.} The server decides, at the moment of sending, what leaves: every answer
of a session in Privacy Mode (JSON, files, exports, the VitalStream viewer's answers) passes
one privacy gate, and the last layer of the server refuses any answer the gate did not clear.
The page shows only what it is sent.

\begin{phibox}
The PHI permission, the roles and Privacy Mode's settings are described in
\cref{chap:users-privacy}. For a review before a deployment on patient data, \cref{part:design}
gathers the architecture (\cref{chap:architecture}), the protection of patient data
(\cref{chap:protecting}) and the security review, with its standards, its deployment checklist
and its limits (\cref{chap:security}).
\end{phibox}

\section{The synthetic demo and examples}
\label{sec:introduction-demo}

Every picture, value, signal and patient of this manual is synthetic. They come from two
sources, neither of which holds any patient data:

\begin{itemize}
\item \textbf{The demo}, made by \file{tools/make-demo-data.py}: six invented patients
  (\code{P900001} to \code{P900006}, with the MRNs \code{SYN0000001} to \code{SYN0000006})
  with their encounters and bed stays, laboratory results, vital signs, medications,
  diagnoses, microbiology, allergies and notes; short ICU recordings, a 24-hour Holter, an
  overnight EDF+ sleep study, 12-lead ECGs (WFDB, GE MUSE XML, DICOM, PDF and PNG), chest X-rays and CT series, echo loops
  and reports, a coronary angiography run, and another hospital's OMOP database that adds
  earlier stays and results to two of them. The installers make it when asked (\opt{--demo}),
  and \cmd{medorama init} loads it into a cohort named \code{DEMO} that every user sees
  (\cref{chap:quickstart}). It is off by default.
\item \textbf{The viewers' examples}: \file{viewers/ecg/examples/} (a WFDB record, a GE MUSE
  XML file, a DICOM ECG, a PDF and a PNG printout) and \file{viewers/pacs/examples/} (an echo
  loop, a coronary angiography run, a chest X-ray), written by
  \file{tools/make-viewer-examples.py}; \file{viewers/vitalstream/testdata/} (a WFDB
  collection, plain records, a folder of records to open as one session, and EDF, EDF+ and BDF+
  files).
\end{itemize}

The repository itself holds no patient data: the folders \file{sample\_data/} and
\file{viewers/vitalstream/sample\_data/} are for local data and are never committed, and a
privacy check, run by the git hooks before each commit and by CI on every commit, refuses
recordings, data files and real recording names.

\section{The website}
\label{sec:introduction-website}

The website, \url{https://medorama.org}, holds the documentation, the installation commands
for each operating system (\url{https://medorama.org/docs/install.html}), and
runs the ECG and PACS viewers in the browser, on the synthetic examples or on the visitor's
own files, which are read on their own computer and never leave it. VitalStream is a server
and does not run on the website. The source code is at
\url{https://github.com/alphanumericslab/medorama}.

\section{This manual}
\label{sec:introduction-manual}

The manual is written mainly for users, administrators and integrators, but it covers every
feature, the API and the installation. It gathers the documents of the repository
(\file{docs/}, the READMEs) into one book, and agrees with them and with the code. It has seven
parts and appendices:
\begin{itemize}
\item \cref{part:start}: this introduction, the installation (\cref{chap:installation}) and a
  first tour of the demo (\cref{chap:quickstart});
\item \cref{part:using}: the page, screen by screen;
\item \cref{part:data}: bringing data in, as imported tables, connected databases, CLIF,
  recordings, ECGs and imaging;
\item \cref{part:viewers}: the viewers on their own;
\item \cref{part:admin}: administration, users and privacy, sign-in, deployment, add-ons,
  readers and new formats, troubleshooting;
\item \cref{part:design}: the architecture, the protection of patient data and the security
  review, written so that an institution's IT security office can review the design on its
  own;
\item \cref{part:reference}: the command line, the HTTP APIs, the settings and the formats;
\item the appendices: a glossary (\cref{chap:glossary}) and the licenses
  (\cref{chap:licences}).
\end{itemize}

\subsection{Conventions}
\label{sec:introduction-conventions}

\Cref{tab:introduction-markup} shows how the text marks what it names.

\begin{table}[htbp]
\caption{The conventions of this manual.}
\label{tab:introduction-markup}
\begin{tabularx}{\linewidth}{@{}P{4.2cm}L@{}}
\toprule
{\tabhead Written as} & {\tabhead Means} \\
\midrule
\ui{Patient Catalog}, \ui{Add to cohort\dots} & A tab, button or menu of a page, written
  exactly as the page shows it \\
\key{Enter}, \key{Esc} & A key of the keyboard \\
\code{QTc > 470ms}, \code{results\_dir} & A value, an identifier, a setting, an expression \\
\cmd{medorama serve} & A command line or a subcommand \\
\opt{-{}-demo} & A command-line option \\
\file{tools/install-medorama.sh} & A file or folder, of the repository or of an output \\
name\phimark & Protected health information (an identifier) \\
\bottomrule
\end{tabularx}
\end{table}

Command lines are shown in listings like this one:
\begin{lstlisting}[style=shell]
medorama --version
\end{lstlisting}
and what a program prints, when it matters, in a listing of its own:
\begin{lstlisting}[style=text]
medorama 0.1.0
\end{lstlisting}

Three kinds of boxes set text apart:
\begin{note}
A note: something to know before going on.
\end{note}
\begin{tip}
A tip: a quicker or better way.
\end{tip}
\begin{phibox}
Patient data: what concerns identifiers and Privacy Mode.
\end{phibox}

The screenshots show the dark theme, except those of the ECG viewer, which is shown in both
themes (\cref{chap:ecg-viewer}). Paths are written with \code{/}; on Windows, read \code{\textbackslash}.
Commands for Linux and macOS are given first, and their Windows equivalents where they differ.

\section{How to cite Medorama}
\label{sec:introduction-cite}

The file \file{CITATION.cff} at the root of the repository says how to cite Medorama (or
VitalStream): Reza Sameni, \emph{Medorama: A unified patient-journey research platform
linking EHR data, physiological waveforms, imaging and laboratory results across databases},
version 0.1.0, released 2026-10-11, license PolyForm Noncommercial 1.0.0,
\url{https://github.com/alphanumericslab/medorama}. In BibTeX:
\begin{lstlisting}[style=text]
@software{sameni_medorama,
  author  = {Sameni, Reza},
  title   = {Medorama: A unified patient-journey research platform linking EHR data,
             physiological waveforms, imaging and laboratory results across databases},
  version = {0.1.0},
  date    = {2026-10-11},
  license = {PolyForm-Noncommercial-1.0.0},
  url     = {https://github.com/alphanumericslab/medorama}
}
\end{lstlisting}

\paragraph{License.} Medorama is licensed under the PolyForm Noncommercial License 1.0.0: free
for noncommercial use, such as research and teaching and use by universities, public research
and public health organizations; commercial use needs a commercial license from the author. WFDB
records are read with wfdb-x, a separate project under the GNU GPL version 3 or later; every
program can also be built without WFDB support, and then contains no wfdb-x code
(\cref{sec:installation-no-wfdb}). Every dependency, with its license, is listed in
\cref{chap:licences}. Product and company names used in this
manual (GE, MUSE, DICOM, OMOP and others) are trademarks of their respective owners; they only
name the systems and formats Medorama reads or connects to.
